\subsection{Health and safety impacts}

Health is also a critical dimension in \acp{bss} promotion, complementing the environmental and transport outcomes discussed earlier~\cite{DeMaio2009, Fishman2016, Qiu2018}. This section explores the health and safety implications of bike-sharing, beginning with evidence of user-level physical and mental benefits, then reviewing the results of \acp{hia} studies that quantify population-level effects through physical activity, air pollution, and traffic-safety pathways, and concluding with a discussion of equity considerations and policy implications.

    \subsubsection{User-level health benefits of cycling}

    Regular cycling has been consistently associated with substantial health benefits, including improved cardiovascular fitness and lower risks of diabetes, cancer, and premature mortality~\cite{Green2021, Celis-Morales2017, Dutheil2020}. It provides an efficient means of achieving the recommended 150 minutes of moderate to vigorous physical activity per week~\cite{Knowler2002}, while also supporting mental health by reducing stress, alleviating symptoms of depression, and enhancing overall well-being~\cite{Chekroud2018, Ruiz2015}. By extending access to bicycles and encouraging regular use, \acp{bss} have the potential to translate these individual-level health benefits into broader population gains, but to what extent have such impacts been empirically demonstrated?

    Over the past decade, a growing body of research has examined how bike-sharing affects public health, particularly through changes in physical activity. Early observational and survey-based studies provided the first evidence of user-level health improvements. In 2010, Montreal's \textit{BIXI} system data showed that residents living near docking stations increased their cycling activity, suggesting that \acp{bss} can promote active travel~\cite{Fuller2013}. However, the study did not determine whether this increase came at the expense of other forms of active mobility. Around the same period, an eight-month study among university commuters in Valencia found that regular use of \textit{Valenbisi} contributed roughly half of the recommended 150 minutes of weekly moderate physical activity and was associated with modest reductions in BMI~\cite{Molina-Garcia2013}. A survey of 3,000 \textit{Capital Bikeshare} users in Washington, D.C., similarly found that 31.5\% reported reduced stress and 30\% reported weight loss as a result of using the system~\cite{Alberts2012}. Complementing these self-reported outcomes, Xu~\cite{Xu2019} linked U.S. bike-sharing deployment data with county-level obesity prevalence statistics and found an average 1\% decline in obesity rates following \acp{bss} implementation~\cite{Hosford2019}.

    Expanding to a multi-city perspective, Fishman et al.~\cite{Fishman2014} analyzed systems in Melbourne, Brisbane, Washington D.C., London, and Minneapolis–St. Paul, reporting substantial overall increases in physical activity (from about 1.4 million additional minutes per year in Minneapolis–St. Paul to 74 million in London). However, the authors emphasized that the magnitude of health gains depends strongly on modal substitution: when bike-sharing trips replace car travel, users gain additional physical activity and related health improvements, whereas substituting walking trips produces the opposite effects. Similar results were reported for \textit{Dublinbikes}, where many trips replaced walking rather than motorized travel, leading to modest health improvements~\cite{Murphy2015}. Therefore, health benefits are contingent not only on \ac{bss} usage levels but also on the modes replaced.

    Taken together, these studies demonstrate that bike-sharing can effectively promote cycling and deliver measurable improvements in individual fitness and mental well-being. Yet translating these user-level benefits into citywide health gains remains challenging. A narrative review by Bauman et~al.~\cite{Bauman2017} found little evidence that current systems substantially increase overall physical activity for the general population, as most \ac{bss} users were already active and many trips substituted for walking or \ac{pt} rather than car travel. According to the authors, because regular users account for only a small share of the urban population, even if all met \ac{who} activity guidelines through cycling, the proportion of adults achieving these recommendations would rise by less than one percentage point. This suggests that significant health improvements at a city level are likely to occur only when bike-sharing reaches much broader segments of society.
        

    \subsubsection[HIA framework]{\ac{hia} framework}

    While the previous empirical studies provide valuable insights into individual effects, they offer limited evidence of citywide impacts or the specific pathways through which health effects occur. To address this gap, researchers have employed \acp{hia} as modelling frameworks to estimate the combined effects of bike-sharing on health outcomes. \acp{hia} estimate the net population-level health impacts by combining travel behavior data with epidemiological dose–response relationships~\cite{Rojas-Rueda2011, Woodcock2014, Otero2018, Clockston2021}. They typically evaluate changes in key outcomes such as mortality, morbidity, and \acp{daly}, accounting for three principal pathways~\cite{Teixeira2021}: (i) health gains from increased physical activity, (ii) health risks from traffic injuries, and (iii) health losses from exposure to air pollution. This holistic approach allows researchers to study both the benefits and risks of \acp{bss} under different scenarios.

    To quantify physical activity, these models use the concept of the \ac{met}, a standardized physiological measure that expresses the energy expenditure of a specific activity as a multiple of the resting metabolic rate~\cite{Ainsworth2011}. One \ac{met} corresponds to the energy cost of sitting quietly, equivalent to an oxygen uptake of approximately 3.5~\si{\milli\liter}~\ce{O2} per kilogram of body weight per minute. Thus, an activity with an intensity of 6~\acp{met} requires roughly six times the energy expended at rest. Cycling intensity is typically assumed to range between 6 and 8~\acp{met} for conventional bicycles and around 6~\acp{met} for e-bikes with standard assistance~\cite{Rojas-Rueda2011, Woodcock2014, Otero2018, Clockston2021}. In this context, the total dose of physical activity is first estimated at the individual level, based on trip duration, frequency, and the proportion of the population using the \ac{bss}. This dose is expressed in \ac{met}-hours per week or per trip, reflecting the additional energy expenditure relative to rest. The resulting values are then combined with epidemiological dose–response functions derived from large cohort studies or meta-analyses to estimate the corresponding reductions in disease risk. Because these relationships are empirically non-linear, some \acp{hia} incorporate functions that assign larger relative benefits to previously sedentary populations and smaller marginal gains to already active individuals~\cite{Otero2018, Clockston2021, Chen2023}. Through this process, \acp{hia} translate individual-level physical activity increases associated with \acp{bss} into population-level estimates of prevented deaths or gains in \acp{daly}.

    Air-pollution exposure is typically assessed in \acp{hia} through the inhalation of fine particulate matter (\ce{PM_{2.5}}) during travel, as this pollutant shows the strongest association with all-cause mortality and disease incidence~\cite{Woodcock2014, Otero2018, Clockston2021, Cao2021}. The inhaled dose is computed by combining background \ce{PM_{2.5}} concentrations with activity-specific ventilation rates, which are higher for cyclists than for pedestrians or motorists, and trip duration. This value represents the total pollutant mass inhaled per trip, which is then linked to health outcomes using the dose–response functions derived from epidemiological studies. Even though other pollutants such as \ce{NO2} or black carbon could also be modelled, most studies focus on \ce{PM_{2.5}} as they are highly correlated and produce similar health responses~\cite{Otero2018}. Moreover, for simplicity, assessments restrict the analysis to in-transit exposure and do not account for the indirect improvements in air quality resulting from reduced car use~\cite{Clockston2021}. 

    For context, the \ac{who} recommends since 2021 long-term average \ce{PM_{2.5}} concentrations not exceeding 5~\si{\micro g/m^3} and daily averages below 15~\si{\micro g/m^3}~\cite{WHO2021}. However, air-quality levels in many cities, particularly in low- and middle-income countries, remain well above these thresholds. In Chinese megacities, for example, annual mean \ce{PM_{2.5}} concentrations on from 2004 to 2013 have ranged from 32 to 120~\si{\micro g/m^3}~\cite{Ma2016}, often exceeding international guidelines by an order of magnitude. Such elevated pollution levels imply that even modest increases in active travel can entail greater inhalation of pollutants, which may offset part of the health benefits from physical activity.

    Finally, traffic-safety impacts are measured by estimating the number of injuries and fatalities expected when substituting other travel modes with cycling. Models combine exposure metrics such as kilometres or hours cycled with mode-specific crash rates derived from police or hospital records~\cite{Otero2018, Clockston2021}. These rates are then applied to the estimated number of new cycling trips to calculate changes in injury and fatality risk relative to the baseline.

    
    \subsubsection[Empirical findings from HIAs]{Empirical findings from \acp{hia}}

    \acp{hia} have consistently found that the overall health impacts of \acp{bss} are driven primarily by gains in physical activity, while the air-pollution and traffic-injury pathways play comparatively minor roles. This section reviews empirical applications across different contexts, beginning with the main findings and the dominant influence of the physical-activity pathway, and then examining the air-pollution and traffic-injury pathways.

    \paragraph{Overview of \acp{hia} applications and the central role of the physical-activity pathway\\}

    One of the first studies of this kind assessed Barcelona’s \textit{Bicing} program~\cite{Rojas-Rueda2011} in 2009. The analysis estimated that increased physical activity among users prevented approximately 12.5 deaths per year, while the combined risks from traffic crashes and pollution exposure were minimal by comparison. The net effect corresponded to 69 deaths avoided per million users annually, yielding a benefit–risk ratio of approximately 77:1. However, this outcome partly reflected an optimistic modeling assumption: the authors assumed that about 90\% of \textit{Bicing} trips replaced car journeys. This strong mode-shift assumption—acknowledged by the authors as a limitation—likely amplified the estimated health benefits. A similar study conducted in 2010 in London’s \textit{Barclays Cycle Hire} incorporated empirical trip data and disaggregated outcomes by sex and age~\cite{Woodcock2014}. Using a comparative risk assessment approach, the authors found a smaller but still positive overall health benefit compared to Barcelona, with the magnitude of gains varying across demographic groups: men experienced larger benefits because they cycled more frequently and for longer distances, older adults benefited more in absolute terms due to higher baseline mortality risk reductions, while women faced higher relative injury risks in traffic.     
    
    Building on these city-level analyses, Otero et al.~\cite{Otero2018} assessed 12 large-scale \acp{bss} across Europe in 2018, estimating a combined 5.17 deaths avoided per 1,000 bicycles in service per year. The study explicitly included both \acp{m-b} and \acp{e-b}, distinguishing between standard- and high-assistance \acp{e-b}. Although \acp{e-b} provide lower levels of physical exertion per trip, the authors argue that they can still have a positive health impact, as they attract a wider range of users. Overall, the authors reported a benefit–risk ratio of 19:1, indicating that physical-activity gains outweighed risks from air pollution and traffic injuries by a wide margin. They also found that larger systems tend to generate greater aggregate health benefits, as the higher number of users and trips amplifies population-level effects. Although absolute numbers remained modest, the study highlighted that benefits could scale up substantially if a higher proportion of \acp{bss} trips replaced car journeys; for instance, increasing the share of car-substituted trips to 50\% could raise the number of deaths avoided nearly tenfold.
    
    In the U.S., Clockston et al.~\cite{Clockston2021} conducted the first nationwide \ac{hia} of \acp{bss}, using 2019 ridership data to quantify both mortality and morbidity effects. Across all U.S. systems, \ac{bss} use was estimated to prevent approximately 4.7 premature deaths and 737~\acp{daly} per year, equivalent to about \$36 million in avoided health costs. The authors also examined New York City separately, where \ac{bss} trips prevented around 2 premature deaths and 355~\acp{daly} annually, yielding health savings of roughly \$15 million. Unlike the earlier European analysis by Otero et al.~\cite{Otero2018}, which focused solely on mortality, this assessment incorporated morbidity through \acp{daly} and economic valuation. The study further distinguished among substitution modes showing that the largest health gains occurred when \ac{bss} replaced car trips, followed by walking, while replacing \ac{pt} trips produced smaller or even neutral effects. Consistent with previous evidence, the increase in physical activity accounted for about 90\% of the total health impacts.

    Lastly, extending this evidence beyond Europe and North America, Chen et al.~\cite{Chen2023} analyzed Shanghai’s \textit{Mobike} program, offering the first \ac{hia} of a \ac{bss} in a highly polluted Asian megacity. Even under average \ce{PM_{2.5}} concentrations of about 45~\si{\micro g/m^3}, the study found that the gains from physical activity clearly dominated, preventing roughly 7–8 premature deaths per year among 306,000 users. Therefore, the overall health balance remained positive up to very high pollution levels. Only when \ce{PM_{2.5}} concentrations approached 68~\si{\micro g/m^3} did the pollution burden begin to offset the benefits. These results reaffirm the central role of physical activity as the main driver of health gains from \acp{bss}, even in environments with substantial air-quality challenges.


    \paragraph{Air-pollution pathway\\}

    Across studies, the air-pollution pathway has been found to contribute only marginally to the overall health balance of \acp{bss}. In both Barcelona and London, the additional exposure to \ce{PM_{2.5}} during cycling was outweighed many times over by the benefits of increased physical activity, with the balance turning negative only under unrealistically high pollution scenarios~\cite{Rojas-Rueda2011, Woodcock2014}. Extending the analysis to 12 European cities, Otero et al.~\cite{Otero2018} similarly identified \ce{PM_{2.5}} as the least influential of the three modelled pathways but noted that none of the cities met the \ac{who}’s air-quality guidelines, suggesting that cleaner urban air would further amplify net benefits.
    
    Consistent with this pattern, Clockston et al.~\cite{Clockston2021} confirmed that pollution effects remain small at the national scale, estimating less than one death attributable to in-transit \ce{PM_{2.5}} exposure per year. Their analysis also showed that the health impact of air pollution depends strongly on which modes are replaced. A counterintuitive result emerged when bike-sharing trips substituted for walking: exposure to \ce{PM_{2.5}} slightly decreased because cycling trips were shorter in duration, leading to lower total inhaled volumes despite higher ventilation rates. By contrast, when trips replaced car or \ac{pt}, exposure increased, as cyclists breathe more intensely and are directly immersed in the traffic stream~\cite{Rojas-Rueda2011, Otero2018, Clockston2021}. Overall, however, the evidence consistently shows that these pollution-related risks remain minor compared with the substantial health benefits derived from increased physical activity.

    Nevertheless, a few studies have shown that the air-pollution pathway can play a more prominent role under specific conditions. Chen et al.~\cite{Chen2023} reported that in Shanghai’s \textit{Mobike} program, elevated \ce{PM_{2.5}} concentrations significantly affected the overall health balance. Under average levels of about 45~\si{\micro g/m^3}, the benefits from physical activity still outweighed the harms from pollution, but as concentrations approached 68~\si{\micro g/m^3}, the air-pollution burden began to offset these gains, eventually reversing the net outcome. This finding highlights that in highly polluted cities, the contribution of the air-pollution pathway can become decisive, potentially turning otherwise beneficial cycling activity into a net health loss.
    
    Conversely, Hou et al.~\cite{Hou2023} found a great positive air-quality impact in their assessment of New York City's \textit{Citi Bike} program. The authors estimated an overall benefit of 5,745~\acp{daly} and 216 premature deaths avoided, with major contributions stemming from reductions in \ce{PM_{2.5}} emissions (avoiding 2,757~\acp{daly} and 164 deaths). The prominence of the air-pollution pathway reflects both methodological and contextual differences: their model accounted for system-wide emission reductions resulting from decreased car use, whereas most earlier \acp{hia} considered only users' in-transit exposure. Moreover, New York's high baseline walking rate of approximately 35\% limited the additional physical-activity gains achievable through cycling.

    Beyond \acp{hia} frameworks and their traditional reliance on daily mean \ce{PM_{2.5}} values, recent studies have also underscored the importance of short-term peaks and behavioural adaptation in shaping cyclists’ real-world exposure. The analysis of more than one million \textit{Mobike} trips in Guangzhou by Cao et al.~\cite{Cao2021} revealed that riders were exposed to markedly higher \ce{PM_{2.5}} risks during evening traffic peaks, when exceedance concentrations clustered around dense commercial and leisure areas. This approach demonstrated that daily averages can underestimate inhaled doses by up to an order of magnitude and underscored the need for high-resolution spatiotemporal modelling in future \acp{hia}. Complementing this study, Park et al.~\cite{Park2022} analysed three years of data from Seoul’s public \ac{bss} and found a clear non-linear response between ambient \ce{PM_{2.5}} and ridership: usage declined sharply once concentrations exceeded roughly 38–42~\si{\micro g/m^3}. Commuting trips showed pronounced risk-averting behaviour, whereas leisure trips remained less sensitive to pollution levels.


    \paragraph{Traffic-safety pathway\\}

    While the potential risks associated with traffic injuries may be a concern in discussions about bike-sharing, evidence consistently shows that the traffic-injury pathway contributes far less to overall health outcomes compared to the substantial gains from increased physical activity. Nevertheless, the impact of this pathway is variable and closely linked to local traffic conditions and infrastructure.

    For example, Otero et al.~\cite{Otero2018} found that, across several European cities, the traffic-injury component remained a minor part of the total health balance. However, they observed notable differences among cities: except in Seville, switching from car to cycling via bike-sharing marginally elevated the risk of traffic fatalities, with risks being highest in places like Milan, Brussels, and Warsaw. Such disparities were attributed primarily to differences in cycling infrastructure and road-safety environments. Despite these variations, losses due to traffic injury represented only a small fraction of the total, with the benefit-to-risk ratio remaining favourable at around 19 to 1, reaffirming that health benefits from physical activity consistently outweigh traffic injury risks.

    This European pattern is echoed in the U.S., where Clockston et al.~\cite{Clockston2021} conducted nationwide modelling. They showed that substituting all modes with bike-sharing marginally increased the average user’s risk of traffic incidents—about 0.26 additional fatalities per 100,000 users per year—whereas replacing walking trips actually reduced risk by approximately 0.15. This nuanced outcome reflects U.S. statistics, where pedestrians experience higher fatality rates per distance than cyclists, so shifting trips from walking to cycling can lower total risk. More broadly, just as in European cases, U.S. modelling found that traffic injuries remained a minor factor compared to the gains from increased physical activity, and that the overall effect was highly sensitive to local safety infrastructure. Supporting these findings, Rojas-Rueda et~al.~\cite{Rojas-Rueda2011} also observed similar trends. In New York City specifically, Hou et al.~\cite{Hou2023} even found that traffic-injury risks could yield measurable safety benefits, as reduced car use through bike-sharing helped to decrease crash rates citywide.

    Further detailing the safety profile of bike-sharing, evidence from London illustrates that bike-sharing users may experience injury rates at least as low as those of private cyclists~\cite{Woodcock2014}. In this \ac{hia}, the lower-than-expected injury rates among bike-sharing users compared to private cyclists led to a net positive health effect. Nevertheless, the study also observed that outcomes varied by gender, with women facing higher fatality risks, especially from collisions with heavy goods vehicles. The authors pointed to several possible explanations for these safer bike-sharing outcomes: built-in bicycle features such as lower speeds and lighting, behavioral patterns like greater cycling in parks, and increased awareness among other road users. Syntheses across North American and European cities generally confirm lower crash rates for bike-sharing compared to private cycling~\cite{Martin2016}. Fishman and Schepers~\cite{Fishman2016_road_safety}, analyzing injury data from both continents, reinforced this point. They found that after the introduction of bike-sharing, overall cycling injury risk declined, with bike-sharing users approximately half as likely as private cyclists to suffer fatal or severe injuries.

    Delving into specific safety mechanisms, several studies have highlighted the importance of helmet use, rider experience, and the overall traffic context. Helmet use, in particular, has remained a persistent concern. A systematic review found high unhelmeted rates among bike-sharing users (typically 36--89\%) and observed mixed evidence relating to injury rates after bike-sharing scheme launches~\cite{Chen2021}. For instance, one multi-city analysis showed that, in jurisdictions without mandatory helmet laws, the proportion of head-injured cyclists rose from 42\% to 50\% following bike-sharing implementation, meaning a 30\% higher adjusted odds of head injury~\cite{Graves2014}. In response, some cities have explored options such as helmet rental or educational initiatives, though these measures may sometimes undermine the spontaneity that makes bike-sharing appealing~\cite{Chen2021}. In addition to helmet use, factors such as rider experience and traffic environment exert considerable influence on safety outcomes. Because bike-sharing systems often attract new cyclists there is a particular need for investment in safe cycling infrastructure like protected lanes and traffic-calmed streets~\cite{Otero2018}. Such provisions not only help reduce individual risk but also support broader, community-level safety.

    Finally, it is important to recognize the broader “safety-in-numbers” phenomenon: as the number of cyclists increases, the injury risk per cyclist typically declines. Multiple reviews identify this effect to help explain why growing usage of bike-sharing programs does not result in proportionate increases in injuries, even as cycling activity rises overall~\cite{Jacobsen2015, Teixeira2021}. Thus, effective integration of bike-sharing into urban mobility, along with investments in safety infrastructure, can help ensure that physical-activity gains are realized with minimal risk.
    
    
    \subsubsection{Equity, context, and methodological considerations}

    Although \acp{bss} deliver clear population-level health benefits, these gains are not necessarily distributed equally across all urban spaces or social groups. For example, in New York City, the historical placement of bike-sharing stations has tended to favour lower-poverty neighbourhoods, effectively shaping who can access and benefit most from the system. As the network expanded, estimates of premature deaths averted increased from two to three per year, yet sensitivity analyses suggested that prioritizing expansion in higher-poverty areas could further amplify these benefits~\cite{Babagoli2019}. This pattern of uneven access and impact is echoed in national-level \acp{hia} from the U.S. and New York, which, despite affirming overall net health gains, underscore the necessity of addressing equity through both better data collection on user socio-economic status and a more balanced spatial distribution of stations~\cite{Clockston2021}. Ultimately, ensuring that health improvements are shared across diverse population groups requires deliberate action, particularly as most studies to date have focused on high-income cities characterized by lower air pollution and stronger road safety regulations. In contrast, contexts with poorer air quality or weaker safety environments may experience a markedly different balance of health benefits and risks, highlighting the importance of broadening the evidence base to include a wider range of urban settings~\cite{Teixeira2021, Tainio2016}.

    Despite methodological advancements, bike-sharing \acp{hia} still face several challenges. Results are highly sensitive to assumptions about which modes bike-sharing trips replace. Because most studies lack precise empirical data on mode substitution, they rely on survey-based or modelled estimates, leading to wide variation in health outcomes. Furthermore, many evaluations focus on short-term or immediate effects, leaving open questions about how best to capture long-term behaviour change—such as whether \acp{bss} foster lasting shifts to regular cycling—which could significantly boost the estimated health gains~\cite{Bauman2017}. Finally, while most research centres on physical activity, there is a notable lack of attention to mental-health outcomes, potentially underestimating the full spectrum of benefits and risks associated with bike sharing~\cite{Bauman2017}.

% \textcolor{red}{Teixeira2022: The strengths and weaknesses of bike sharing as an alternative mode during disruptive public health crisis: A qualitative analysis on the users' motivations during COVID-19 
% \begin{itemize}
%     \item Even during the COVID-19 pandemic, riders noted that using bike-sharing saved time and cost on trips while being a “pleasant” experience. 
%     \item Likewise, qualitative interviews in Lisbon during COVID-19 revealed that users appreciate the exercise and eco-friendly aspects of cycling, explicitly associating bike-sharing use with personal health gains and lower carbon footprint. 
% \end{itemize}}